%!TEX encoding = UTF8
%!TEX root = 0-notes.tex

\chapter*{Préface}


\section*{Syllabus}

\subsection*{Description}

Cette UE est une introduction à l'algèbre linéaire (formalisée au S2) qui se base sur l'intuition issue de la géométrie du plan et de l'espace.
Cela inclut une introduction au calcul matriciel.
L'UE introduit aussi le langage de base des polynômes.

\subsection*{Objectifs}

\begin{enumerate}
	\item 
	Géométrie du plan et de l'espace
	\begin{enumerate}
		\item
		Points, vecteurs, translation par un vecteur, combinaisons linéaires, colinéarité, indépendance, bases, repères et coordonnées, changement de repère, barycentres.
		\item
		Droites et plan (sans coordonnées puis avec), positions relatives, intersections, équations.
		\item
		Transformations linéaires et affines classiques : homothéties, translations, symétries, projections du plan et de l'espace.
		\item
		Incursion en géométrie euclidienne : produit scalaire, orthogonalité, distance, produit vectoriel, bases et repères orthonormés, projections orthogonales, distance d'un point à une droite, un plan.
	\end{enumerate}
	\item 
	Algèbre linéaire dans $\R^2, \R^3, $ et $\R^n$ 
	\begin{enumerate}
		\item
		Points et vecteurs de $\R^n$, sous-espaces affines et sous-espaces vectoriels de $\R^n$, expression paramétrique et équations, sous-espace vectoriel engendré par une famille de vecteurs, sous-espace affine engendré par un point et un sous-espace vectoriel.
		\item
		Systèmes linéaires et méthode du pivot : systèmes, ensembles de solutions, matrice d'un système, systèmes échelonnés réduits, opérations élémentaires sur les lignes.
		\item
		Calcul matriciel : opération sur les matrices, matrices des opérations élémentaires sur les lignes.
		\item
		Applications linéaires de $\R^2, \R^3,$ et $\R^n$.
		\item
		Inversibilité d'une matrice et méthode de Gauss-Jordan.
	\end{enumerate}
	\item 
	Polynômes à coefficients réels.
	\begin{enumerate}
		\item
		Définitions d'un polynôme et d'une fonction polynomiale, liens.
		\item
		Coefficients, degré, racines, opérations.
		\item
		Factorisation et division euclidienne de polynômes.
		\item
		Multiplicité des racines, lien à la dérivée, formule de Taylor pour les polynômes.
	\end{enumerate}
\end{enumerate}

\subsection*{Pré-requis nécessaires}

Programme de mathématiques du lycée (notamment géométrie du plan et de l'espace, et résolution d'équations), a minima la spécialité de première et spécialité mathématiques en terminale ou option mathématiques complémentaires.

\subsection*{Pré-requis recommandés}

Programme de mathématiques du lycée (notamment géométrie du plan et de l'espace, et résolution d'équations), idéalement spécialité mathématiques, voire option mathématiques expertes.

\section*{Structure du cours}

Construction d'un cours de mathématiques : Exemple -- Remarque -- Définition -- Lemme -- Proposition -- Théorème -- Démonstration -- Corollaire.

Voir \emph{Mechanica, volume 1}, Leonhard Euler, 1736 (\href{https://scholarlycommons.pacific.edu/euler-works/15/}{E15}), pour un exemple de cours de mathématiques en latin.

Wolframalpha permet de faire du calcul exact.

\begin{definition*}
Associe un mot de vocabulaire à un objet mathématique.
Un mot de vocabulaire peut avoir un sens différent en mathématiques qu'en langage vernaculaire.
C'est par exemple le cas de « ensemble », « couple », « rationnel », « réel », ...
\end{definition*}

\begin{lemme*}
	Résultat préliminaire, parfois technique, utile à la démonstration d'une proposition plus importante.
\end{lemme*}

\begin{proposition*}
	Proposition importante.
\end{proposition*}

\begin{theorem*}
	Résultat fondamental.
\end{theorem*}

\begin{corollaire*}
	Résultat particulier mais important, découlant quasi directement d'un théorème ou d'une proposition.
\end{corollaire*}

%\newcounter{preface}
\begin{Exercise}[label=1]%[counter=preface]
	Exemple de problème d'application directe.
\end{Exercise}
\begin{Exercise}[difficulty=1, label=2]%, counter=preface]
	Exemple de problème plus théorique.
\end{Exercise}
\begin{Exercise}[difficulty=2, label=3]%, counter=preface]
	Exemple de problème d'approfondissement.
\end{Exercise}
\shipoutExercise
\begin{Answer}[ref=1]
	Solution 1.
\end{Answer}
\begin{Answer}[ref=2]
	Solution 2.
\end{Answer}
\begin{Answer}[ref=3]
	Solution 3.
\end{Answer}
\begin{multicols}{3}
\shipoutAnswer
\end{multicols}
